\subsection{VI)}

{\large\begin{align*}\displaystyle F(w, x, y, z) = \sum(2, 3, 12, 13, 14, 15) \quad x = 6; 7; 10; 11\end{align*}}

\begin{center}
\begin{tabular}{c|cccc|c}
\rowcolor[HTML]{FFCCC9} 
\cellcolor[HTML]{9AFF99}$m$ & $w$ & $x$ & $y$ & $z$ & \cellcolor[HTML]{FFFFC7}$F(w,x,y,z)$ \\ \hline
0                           & 0   & 0   & 0   & 0   &                                      \\
1                           & 0   & 0   & 0   & 1   &                                      \\
2                           & 0   & 0   & 1   & 0   & 1                                    \\
3                           & 0   & 0   & 1   & 1   & 1                                    \\
4                           & 0   & 1   & 0   & 0   &                                      \\
5                           & 0   & 1   & 0   & 1   &                                      \\
6                           & 0   & 1   & 1   & 0   & $x$                                  \\
7                           & 0   & 1   & 1   & 1   & $x$                                  \\
8                           & 1   & 0   & 0   & 0   &                                      \\
9                           & 1   & 0   & 0   & 1   &                                      \\
10                          & 1   & 0   & 1   & 0   & $x$                                  \\
11                          & 1   & 0   & 1   & 1   & $x$                                  \\
12                          & 1   & 1   & 0   & 0   & 1                                    \\
13                          & 1   & 1   & 0   & 1   & 1                                    \\
14                          & 1   & 1   & 1   & 0   & 1                                    \\
15                          & 1   & 1   & 1   & 1   & 1                                   
\end{tabular}
\end{center}

\begin{center}
    \begin{karnaugh-map}[4][4][1][$y\qquad z$][$w\qquad x$]
        \minterms{2,3,12,13,14,15}
        \terms{6,7,10,11}{$x$}
        \implicant{3}{10}
        \implicant{12}{14}
    \end{karnaugh-map}
\end{center}

\begin{itemize}
    \item Primer bloque ($m_3$, $m_2$, $m_7$, $m_6$, $m_15$, $m_14$, $m_11$ y $m_10$) permanece constante $y$, va sin negar.
    \item Segundo bloque ($m_12$, $m_13$, $m_15$ y $m_14$) permanecen constantes $w$ y $x$, van sin negar.
    \item La función simplificada queda:
\end{itemize}
{\large
    \begin{align*}
        F(w,x,y,z) &= y + w\cdot{}x
    \end{align*}
}

\subsubsection{Simulación}
\imagen[]{10cm}{./imagenes/ejercicio6Falstad.png}\subsubsection{Descripción en Verilog}\begin{lstlisting}
module descripcionVerilogEjercicio6(
    input inputW,
    input inputX,
    input inputY,
    input inputZ,
    output outputFunction
    );
	 assign outputFunction = inputY | inputW & inputX;


endmodule\end{lstlisting}\subsubsection{RTL output}\imagen[Bloque RTL]{15cm}{./imagenes/rtlEjercicio6Verilog.png}\imagen[Implementación con compuertas]{15cm}{./imagenes/rtlCompuertasEjercicio6Verilog.png}\subsubsection{Simulación temporal}\noindent Código del módulo Verilog Test Fixture:\begin{lstlisting}
module verilogTestFixtureEjercicio6;

	// Inputs
	reg inputW;
	reg inputX;
	reg inputY;
	reg inputZ;

	// Outputs
	wire outputFunction;

	// Instantiate the Unit Under Test (UUT)
	descripcionVerilogEjercicio6 uut (
		.inputW(inputW), 
		.inputX(inputX), 
		.inputY(inputY), 
		.inputZ(inputZ), 
		.outputFunction(outputFunction)
	);

	initial begin
		// Initialize Inputs
		inputW = 0;
		inputX = 0;
		inputY = 0;
		inputZ = 0;

		// Wait 100 ns for global reset to finish
		#100;
        
		// Add stimulus here
		inputW = 0; inputX = 0; inputY = 0; inputZ = 1;
		#100;
		inputW = 0; inputX = 0; inputY = 1; inputZ = 0;
		#100;
		inputW = 0; inputX = 0; inputY = 1; inputZ = 1;
		#100;
		inputW = 0; inputX = 1; inputY = 0; inputZ = 0;
		#100;
		inputW = 0; inputX = 1; inputY = 0; inputZ = 1;
		#100;
		inputW = 0; inputX = 1; inputY = 1; inputZ = 0;
		#100;
		inputW = 0; inputX = 1; inputY = 1; inputZ = 1;
		#100;
		inputW = 1; inputX = 0; inputY = 0; inputZ = 0;
		#100;
		inputW = 1; inputX = 0; inputY = 0; inputZ = 1;
		#100;
		inputW = 1; inputX = 0; inputY = 1; inputZ = 0;
		#100;
		inputW = 1; inputX = 0; inputY = 1; inputZ = 1;
		#100;
		inputW = 1; inputX = 1; inputY = 0; inputZ = 0;
		#100;
		inputW = 1; inputX = 1; inputY = 0; inputZ = 1;
		#100;
		inputW = 1; inputX = 1; inputY = 1; inputZ = 0;
		#100;
		inputW = 1; inputX = 1; inputY = 1; inputZ = 1;
		#100;
	end
      
endmodule
 \end{lstlisting}\imagen[Simulación temporal]{15cm}{./imagenes/simulacionTemporalEjercicio6.png}
